\subsection{Frontend mit Vite und TypeScript}
\label{subsec:frontend-vite-typescript}

Das Frontend bildet die gemeinsame Präsentationsschicht aller fünf Profile.
Vite stellt im Entwicklungsbetrieb den Modulserver bereit und erzeugt mit
\texttt{vite build} statische Artefakte; der vorgeschaltete TypeScript-Lauf
prüft den strikt konfigurierten Quellcode ohne JavaScript-Ausgabe. Vitest deckt
den HTTP-Client ab. Diese Aufgabenteilung entspricht den dokumentierten Rollen
von Vite als Entwicklungs- und Build-Werkzeug sowie TypeScript als statischer
Typprüfung \parencite{vite2026guide,typescript2026handbook}.

\path{frontend/src/main.ts} liest das generierte Profilmodul und zeigt die
Backend-Integration nur bei aktiviertem Feature. Alle Aufrufe liegen in
\path{frontend/src/api/}; \texttt{backend.ts} verlangt dann eine öffentliche
\texttt{VITE\_API\_BASE\_URL} und ruft \texttt{/api/health} auf. Dieselben
Quellen werden im Browser und in Tauri verwendet. Serverlogik, Geheimnisse und
direkter Datenbankzugriff bleiben außerhalb dieser Grenze; Web-only und
Desktop-local benötigen folglich kein Backend
\parencite{kleiveist2026architecture,kleiveist2026profiles}.

Die im untersuchten Commit deklarierten Bereiche und aufgelösten Versionen sind
in Tabelle~\ref{tab:technology-stack} zusammengefasst.

\begin{table}[H]
  \centering
  \small
  \renewcommand{\arraystretch}{1.2}
  \caption{Struktur zur Dokumentation des Technologie-Stacks}
  \label{tab:technology-stack}
  \begin{tabularx}{\textwidth}{@{}p{0.20\textwidth}p{0.27\textwidth}p{0.20\textwidth}>{\raggedright\arraybackslash}X@{}}
    \toprule
    \textbf{Technologie} & \textbf{Rolle im Template} & \textbf{Version / Quelle} & \textbf{Projektbeleg} \\
    \midrule
    % Später ausschließlich tatsächlich verwendete Technologien eintragen.
    % Versionen, Rollen und Belege aus Projektartefakten und Primärquellen ableiten.
    \bottomrule
  \end{tabularx}
\end{table}

